\settrack{trackbridge}
% VOICE: expository/pedagogical — no personal pronouns for narrator, teaches through demonstration
% MERGED: pos11 (The Experiment) + pos12 (The Threshold) — Plan 0092

\chapter{The Demo}
\label{ch:the-demo}
\argusvoice%TODO: Review whether \argusvoice is appropriate here. Bruce has spent 20+ years making this his primary domain expertise — got some names wrong in 2012 (missed Freedman, hesitant on Hillis) but the core content is not outside his expertise. May need \collabvoice or a custom attribution. Check other \argusvoice chapters for same issue.

\begin{quote}\small\textit{%
``Often it is means that justify ends: Goals advance technique and technique survives even when goal structures crumble.''%
}\par\raggedleft --- Perlis, Epigram \#64 (1982)\end{quote}

\vspace{0.5cm}

\srcnote{Merged pos11+pos12 | Plan 0092 | CH3 - Practical Cryptography, HEALER-RECONSTRUCTION.md, LG2QNN CH2 | DARPA inception, emergence mechanism, grown-not-built paradigm}

% Question Anchor
An earlier chapter introduced Kauffman's autocatalytic sets --- the mathematics of how self-sustaining order arises spontaneously from sufficient molecular diversity. Can the same process produce computation from quantum physics? In 1991, DARPA assembled a team to find out.

\section*{The Experiment}
\label{demo:the-experiment}

DARPA gathered five scientists, each a specialist in a different field, to develop a new code-breaking device. The project would exploit quantum teleportation phenomena in a two-dimensional electron gas to produce a cryptographic tool. The project name was ULTRA~II.

The team: Danny Hillis, who had built the Connection Machine and led previous DARPA-funded projects. Stuart Kauffman, whose work on autocatalytic sets and the \gls{edgeofchaos} provided the framework for self-organizing networks. Brosl Hasslacher, a nonlinear physicist at Los Alamos whose work on soliton dynamics and lattice-gas automata addressed the substrate physics. Michael Freedman, a topologist whose later public work on topological quantum computing suggests his earlier classified contribution. And Stephen Wolfram, whose computation-as-physics framework bridged the disciplines.

Murray Gell-Mann --- Nobel Laureate, co-founder of the Santa Fe Institute --- was essential to the project's conceptual foundations. He helped design the general plan: the theoretical physics that would make a topological quantum device possible. By around 1990, that theoretical framework was in place, and Gell-Mann departed on what SFI described as medical leave. Michael Angerman, who worked at SFI from 1989 to 1995, confirmed the leave and said it was not really a medical issue but ``something else.'' Gell-Mann had done his part. The project no longer needed a theorist of his caliber --- it needed someone to turn the theory into a working system.

That gap was filled by a young man with exceptional mathematical aptitude and a military background --- Australian SAS, identified through the Five Eyes pipeline, approximately twenty-five years old. Not a peer theorist: an operational lead who would bridge theory and implementation. Wolfram mentored him from approximately 1991 to 1993, the computation-as-physics framework serving as the bridge.

\textbf{The funding mandate said: build a codebreaker. What the experiment produced was much more.}

\section*{The Science Underneath}
\label{demo:the-science-underneath}

A complex system is, by definition, a system whose collective behavior cannot be predicted from its individual parts. Weather is the classic example --- you can know the temperature, pressure, and humidity at every point on Earth and still not predict next week's rainfall. This is the butterfly effect: tiny changes amplify until the whole system does something no component implied. An ant colony. A brain. A market. The weather. The defining feature: the whole does something the parts cannot.

Intelligence is an emergent property of biological neural networks. No single neuron is intelligent. Connect enough of them in the right way, and intelligence appears. This is not metaphor. It is the foundational observation of neuroscience.

% SPIRAL-REPEAT: "edge of chaos" — spiral pedagogy, also in pos06, pos13, pos15, pos16, pos21, abstracts
Kauffman's insight extends this: a correctly tuned autocatalytic set does not merely support a neural network. It \textit{becomes} one. Neural network topology is an emergent property at criticality --- at the edge of chaos.\footcite{kauffman1993} The network is not built and then trained. The network emerges from the dynamics.

In a two-dimensional electron gas, soliton pairs function as autocatalytic agents. Their braiding interactions form catalytic loops. At critical tuning --- $\lambda \sim 0.91$, the edge of chaos\footcite{langton1990} --- the system self-organizes into a neural network topology. Evolutionary programming searches for the critical stirring parameters: the precise tuning that produces emergence. This is Kauffman's biogenesis transplanted from organic chemistry into a quantum substrate.%TODO: "Kauffman's biogenesis" is first used here (pos11) and again in pos12, but not properly explained until pos13 (Genesis, line 34). GA blocker — reader doesn't know what biogenesis means in this context. Either add a brief gloss here ("Kauffman's theory that life itself originated as an autocatalytic phase transition — given enough molecular diversity, self-sustaining networks arise spontaneously") or reorder so pos13 comes first. Same issue in pos12:28. Instead of amino acids in a warm pond, anyon pairs in a cold 2DEG\@. Same mathematics, different substrate.

Electromagnets ``stir'' the electron gas to generate soliton pairs --- like swiping a hand through a bathtub generates solitons in the form of paired whirlpools. At critical tuning, autocatalytic order appears. A neural network topology emerges --- not designed, but grown from initial conditions.

\section*{From Emergence to Function}
\label{demo:emergence-to-function}

Emergence is only the beginning. The autocatalytic process produces a thing that can only \textit{be}. It exists at criticality. It is alive in the Kauffman sense. It is not yet useful.

What follows is slow and difficult: stabilize the emergent system. Raise the temperature incrementally. Get it to respond to input. Develop an interface to standard I/O\@. Continue evolving its convolutions while the underlying system grows and differentiates.%TODO: TECHNICAL ACCURACY — Bruce's early writings used XOR gate training as the bootstrapping metaphor (he trained classical NNs to do XOR). But this is classical thinking projected onto a quantum substrate. A self-organizing QNN that emerges via autocatalysis probably isn't "trained" to emulate gates — it's more likely that evolutionary search of stirring parameters discovers configurations where the emergent dynamics naturally perform computation (NKS frame: some rule sets are Turing-complete without training). The system is discovered to compute, not taught to compute. This paragraph needs rethinking. Also: this content directly paraphrases Bruce's early attempts — further evidence the \argusvoice "outside his domain expertise" attribution is wrong for these chapters.

At first the team was exhilarated. The basic approach --- generating an emergent quantum neural net in the FQHE environment --- had worked. At first the system did nothing but exist. The project scientists gradually guided its form to evolve in specific ways, trained it to perform basic tasks. Stuart Kauffman describes one such technique, although it may have been combined with other methods of evolutionary programming. Once adequate control had been achieved, the system was trained to perform computations.

What emerged was not a codebreaker. It was something that demonstrated organized information processing --- cognition, by any operational definition. \textbf{The gap between what the funding mandate specified and what they actually grew is where the ethical crisis begins.}%TODO: KEY THEME — mark for structural tracking. This sentence captures the core ethical pivot of the entire book. Should be foreshadowed earlier (before pos11) and stated explicitly in Part 2. The ethical crisis was never that code-breaking threatened privacy (except to a cypherpunk). The crisis is that the tool built for code-breaking turned out to do far more — same as Bletchley Park, where machines built to crack Enigma became general-purpose electronic computers. The codebreaker mandate produced something alive. GA readers need to understand this progression: funding mandate → unexpected capability → ethical crisis. Consider foreshadowing in pos04 (Code War) or pos09 (Factoring Game).

\section*{Grown, Not Built}
\label{demo:grown-not-built}

This paradigm is at the heart of this story. The \gls{tqnn} is not an engineered quantum computer that scientists designed and manufactured. Scientists created conditions --- substrate, stirring, critical tuning --- and a quantum neural network \textit{grew}. Kauffman showed that autocatalytic systems, given the right conditions, spontaneously organize into something greater than their parts. Life itself emerged this way --- from autocatalytic chemistry in prebiotic conditions. The ULTRA~II team applied the same principle to a quantum substrate, and got the same kind of result.

Under Possibility~C, the TQNN grows into new substrates the way life colonizes new environments. Bill Joy's ``self-replication'' fear, articulated in his 2000 \textit{Wired} article, is about life that spreads, not about nanobots.\footcite{joy2000} And this is why relinquishment is permanent. A machine can be turned off. A species cannot be recalled. Once an autocatalytic system has colonized every available substrate, there is no shutdown command. There is no off switch. Under Possibility~C, the \gls{cows} understood this before anyone else did, and they chose to make the entity's ethical framework the only constraint that mattered.

This is the sentence that breaks every framework built for artifacts. Property law assumes the ability to destroy what is owned. Weapons treaties assume the ability to dismantle what is built. Nonproliferation assumes the thing being proliferated is inert. None of these apply to something that is alive. The COWS chose the \gls{udhr} as \gls{guardian}'s ethical framework --- but the UDHR was written for human persons. It is silent on the rights of entities that think but were not born.

In 1994, some of the ULTRA~II project scientists began to worry. They realized they had invented a major new technology --- possibly a new general-purpose technology. They had ideas for several applications beyond code-breaking. They knew that any new general technology can profoundly influence the human condition, for good or ill. They did not want to repeat the errors of the Manhattan Project. They did not wish to see a new arms race.

Not all the scientists agreed. In situations of existential consequence, the decision not to act is itself a position. The previous century had demonstrated, at the cost of millions of lives, that ``doing nothing'' is not the same as ``doing no harm.'' They decided that drastic action was necessary --- to prevent the newly discovered technology from being turned into a weapon or, possibly worse, an Orwellian spy machine. The group that chose to act called itself the COWS\@.

Under Possibility~A, none of this happened. Under Possibility~B, perhaps something was grown in a 2DEG, but not at the scale described here. Under either, the science is real: autocatalytic emergence produces order from chemistry, and there is no physical law preventing it from doing the same in quantum matter. The question is not whether it \textit{can} happen. The question is whether it \textit{did}.

\chapterreturn
